\section{System Design}\label{sec:design}

\subsection{Components and trust boundaries}
\Cref{fig:arch} shows the components and which side of the trust boundary each
one is on. The victim's PC (browser, and the MitB) is untrusted; the phone
authenticator and the bank server are trusted. The cryptographic core
(\texttt{tba\_core}) is shared by the server and, through cross-checked test
vectors, by the phone; the evaluation harness drives the same core and the
server in process.

\begin{figure}[htbp]
\centering
\begin{tikzpicture}[
  font=\small,
  box/.style={draw,rounded corners,align=center,minimum height=8mm,inner sep=3pt},
  trusted/.style={box,fill=green!8},
  untrusted/.style={box,fill=red!8},
  neutral/.style={box,fill=blue!6},
  >=Latex]
  % untrusted PC
  \node[untrusted] (br) {Browser\\(bank pages)};
  \node[untrusted,below=6mm of br] (mitb) {MitB\\\texttt{mitb\_addon.py}};
  % server
  \node[trusted,right=40mm of br] (web) {\texttt{server/app.py}\\routes + templates};
  \node[trusted,below=6mm of web] (ver) {\texttt{verify.py},\\\texttt{transfers.py},\\\texttt{auth.py}};
  \node[trusted,below=6mm of ver] (db) {SQLite};
  % phone
  \node[trusted,above=12mm of web] (app) {Phone authenticator\\Web Crypto + IndexedDB};
  % core / eval
  \node[neutral,right=22mm of ver] (core) {\texttt{tba\_core}\\encoding, totp,\\hmac\_otp, signature};
  \node[neutral,below=6mm of core] (eval) {\texttt{evaluation}\\harness, VirtualPhone,\\scenarios, bench};
  % edges
  \draw[->] (br) -- node[left,font=\scriptsize]{HTTPS} (mitb);
  \draw[->] (mitb.east) -- node[below,font=\scriptsize]{HTTPS} (web.west|-mitb.east);
  \draw[->] (app) -- node[right,font=\scriptsize]{HTTPS} (web);
  \draw[->,dashed] (app.west) to[bend right=20] node[left,font=\scriptsize]{QR on screen} (br.north);
  \draw[->] (web) -- (ver);
  \draw[->] (ver) -- (db);
  \draw[->] (ver) -- (core);
  \draw[->] (eval) -- (core);
  \draw[->] (eval.north) to[bend right=12] node[right,font=\scriptsize]{TestClient} (web.south);
\end{tikzpicture}
\caption{Architecture and trust boundaries (F1). Green: trusted; red: untrusted;
blue: test/verification code. The QR code is the only channel from the
untrusted browser to the trusted phone, and the phone talks to the server on its
own HTTPS channel.}
\label{fig:arch}
\end{figure}

\subsection{Transaction encoding}
A transaction is a fixed JSON object; every field is validated identically in
Python and in JavaScript (\cref{tab:txfields}). The \emph{canonical} form is
\texttt{json.dumps(tx, sort\_keys=True, separators=(",",":"), ensure\_ascii=True)}
encoded as ASCII. Parsing is strict: duplicate keys, floating-point numbers,
\texttt{NaN}/\texttt{Infinity}, non-ASCII bytes, and any string that does not
re-serialise to the identical bytes are rejected. The transaction QR payload is
\texttt{"tba1:"} followed by the base64url of the canonical bytes; the enrolment
QR payload is \texttt{"tbaenroll1:"} followed by a token. The server always
verifies against the canonical transaction stored in \emph{its own} database,
never a transaction sent by the client.

\begin{table}[htbp]
\centering
\caption{Transaction fields and their constraints.}
\label{tab:txfields}
\small
\begin{tabularx}{\linewidth}{@{}lX@{}}
\toprule
Field & Constraint \\
\midrule
\texttt{v}        & integer, equal to 1 \\
\texttt{tx\_id}   & lower-case UUID \texttt{[0-9a-f]\{8\}-...-[0-9a-f]\{12\}} \\
\texttt{from}, \texttt{to} & \texttt{[0-9]\{6,20\}} \\
\texttt{to\_name} & \texttt{[A-Z0-9 ]\{1,64\}} \\
\texttt{amount}   & integer (not bool), $1 \le \texttt{amount} \le 2^{53}-1$ \\
\texttt{currency} & \texttt{[A-Z]\{3\}} \\
\texttt{nonce}    & 16 random bytes, base64url without padding, exactly 22 chars \\
\texttt{exp}      & integer, equal to creation time $+\,120$ seconds \\
\bottomrule
\end{tabularx}
\end{table}

\subsection{The four modes}
\Cref{tab:crypto} gives, for each mode, what the phone computes and what the
server checks. Each device has its own $K_{\text{totp}}$ (20 bytes) and
$K_{\text{hmac}}$ (32 bytes); the ECDSA key pair is generated on the phone. The
domain-separation tags (\texttt{"TBA-HMAC-v1"}, \texttt{"TBA-SIG-v1"}), joined to
the canonical bytes with a \texttt{0x00} byte, stop a code or signature for one
purpose being reused for another.

\begin{table}[htbp]
\centering
\caption{Cryptography of each mode (T4).}
\label{tab:crypto}
\small
\begin{tabularx}{\linewidth}{@{}lXX@{}}
\toprule
Mode & Phone computes & Server checks \\
\midrule
\texttt{baseline} & $\mathrm{HOTP}(K_{\text{totp}},\lfloor t/30\rfloor)$, HMAC-SHA1, 6 digits & window $\pm1$ step; reject step $\le$ last used step \\
\texttt{hmac} & $\mathrm{DT}(\mathrm{HMAC\text{-}SHA256}(K_{\text{hmac}}, \text{tag}\,\|\,\texttt{0x00}\,\|\,\mathrm{canon}(tx)))$, 8 digits & \texttt{hmac.compare\_digest} \\
\texttt{sig} & $\mathrm{ECDSA\text{-}P256\text{-}SHA256}(sk, \text{tag}\,\|\,\texttt{0x00}\,\|\,\mathrm{canon}(tx))$, raw 64 bytes & convert to DER, verify with public key \\
\texttt{sig\_blind} & same as \texttt{sig} & same as \texttt{sig} \\
\bottomrule
\end{tabularx}
\end{table}

\subsection{Device enrolment}
Enrolment (\cref{fig:enroll}) establishes one device per user. Crucially, once a
device exists the server \emph{refuses} to enrol another from a web session
alone: \texttt{GET /enroll} returns 409 and issues no token. Otherwise an
attacker with only a session (A1, or a stolen password) could enrol their own
key and approve every transaction, in every mode. In the demo, re-enrolment
requires wiping the database (\texttt{make reset}); a real system needs an
out-of-band process (\cref{sec:discussion}).

\begin{figure}[htbp]
\centering
\begin{tikzpicture}[font=\small,>=Latex]
  \def\yl{0}        % top
  \def\yb{-5.4}     % bottom
  \foreach \x/\name in {0/User, 3.2/Browser, 7.4/Server, 11.0/Phone}{
    \node[font=\small\bfseries] at (\x,0.3) {\name};
    \draw[gray] (\x,\yl) -- (\x,\yb);
  }
  \tikzset{msg/.style={->,shorten >=1pt}}
  \draw[msg] (0,-0.5) -- node[above,font=\scriptsize]{login, open /enroll} (3.2,-0.5);
  \draw[msg] (3.2,-1.1) -- node[above,font=\scriptsize]{GET /enroll} (7.4,-1.1);
  \node[align=center,font=\scriptsize,fill=yellow!20,draw,rounded corners] at (7.4,-1.8)
        {device exists $\Rightarrow$ 409, no token};
  \draw[msg] (7.4,-2.5) -- node[above,font=\scriptsize]{QR \texttt{"tbaenroll1:"+token}} (3.2,-2.5);
  \draw[msg,dashed] (3.2,-3.0) -- node[above,font=\scriptsize]{scan QR} (11.0,-3.0);
  \node[align=center,font=\scriptsize,fill=blue!8,draw,rounded corners] at (11.0,-3.6)
        {generate P-256,\\non-extractable};
  \draw[msg] (11.0,-4.2) -- node[above,font=\scriptsize]{POST /api/device/enroll \{token, public\_key\}} (7.4,-4.2);
  \draw[msg] (7.4,-4.9) -- node[above,font=\scriptsize]{\{device\_id, k\_totp, k\_hmac, mode\}} (11.0,-4.9);
\end{tikzpicture}
\caption{Device enrolment (F2). The HMAC keys are sent in the clear over HTTPS
and stored by the phone as non-extractable keys in IndexedDB.}
\label{fig:enroll}
\end{figure}

\subsection{Confirmation under a man-in-the-browser}
\Cref{fig:confirm} shows what happens to a transfer when the browser is
compromised. The MitB rewrites the recipient to the attacker before the request
reaches the server, and masks the returned confirmation page so the browser
still shows the intended recipient. The transaction QR, however, encodes the
transaction the server actually holds, so the phone displays the \emph{true}
recipient. In \texttt{hmac} and \texttt{sig} an attentive user sees the mismatch
and rejects; in \texttt{sig\_blind} the phone shows nothing, so the user cannot
notice. This is the mechanism measured in \cref{sec:eval} and shown in the
screenshots of \cref{fig:mitb}.

\begin{figure}[htbp]
\centering
\begin{tikzpicture}[font=\small,>=Latex]
  \foreach \x/\name in {0/User, 3.0/Browser, 6.0/MitB, 9.0/Server, 12.0/Phone}{
    \node[font=\small\bfseries] at (\x,0.3) {\name};
    \draw[gray] (\x,0) -- (\x,-5.6);
  }
  \tikzset{msg/.style={->,shorten >=1pt}}
  \draw[msg] (0,-0.5) -- node[above,font=\scriptsize]{send 1M to A} (3.0,-0.5);
  \draw[msg] (3.0,-1.0) -- node[above,font=\scriptsize]{POST /transfer to=A} (6.0,-1.0);
  \draw[msg] (6.0,-1.5) -- node[above,font=\scriptsize]{\textbf{to=B (attacker)}} (9.0,-1.5);
  \draw[msg] (9.0,-2.2) -- node[above,font=\scriptsize]{confirm page (tx for B)} (6.0,-2.2);
  \node[align=center,font=\scriptsize,fill=red!10,draw,rounded corners] at (6.0,-2.9)
        {mask HTML:\\show A to user};
  \draw[msg] (6.0,-3.5) -- node[above,font=\scriptsize]{page shows A} (3.0,-3.5);
  \draw[msg,dashed] (3.0,-4.0) -- node[above,font=\scriptsize]{scan QR (tx for B)} (12.0,-4.0);
  \node[align=center,font=\scriptsize,fill=green!10,draw,rounded corners] at (12.0,-4.7)
        {phone shows \textbf{B}\\(the truth)};
  \draw[msg] (12.0,-5.3) -- node[above,font=\scriptsize]{attentive user: Reject (hmac/sig)} (9.0,-5.3);
\end{tikzpicture}
\caption{Confirmation under MitB (F3). The browser is deceived, but the phone
reads the transaction the server holds, so \texttt{hmac}/\texttt{sig} expose the
swap. In \texttt{sig\_blind} the phone shows nothing and the user approves.}
\label{fig:confirm}
\end{figure}
